\subsection{Interface and scope}
Let $s\ge2$, $\lambda=\lceil\log_2s\rceil$, and let $\mathcal L$ be the total
input description length. The network is bipartite and may have parallel edges, has at most
$s$ Boolean edges and maximum degree $\lceil s^{0.9}\rceil$, and has entries in
\[
 (\mathbb Z/2^{P_0}\mathbb Z)[V]/(V^\ell),\qquad P_0,\ell\le s.
\]
Its original entry routines cost at most $\poly(s+\mathcal L)2^{A_0s^{0.9}}$ time
and polynomial space, for fixed interface constants. All local tensor descriptions retain
the source's immutable creation-time representatives.

\begin{theorem}[Relative contraction certificate]\label{thm:85}
Assuming the unchanged source deformation, integrality, recovery and recursive-factorization
identities, the construction below contracts this normalized interface in polynomial space and
\[
 (s+\mathcal L)^{O(1)}
 2^{(C+o(1))s/\sqrt\lambda},\qquad C=84.06627149423657\ldots<85.
\]
In particular $85$ replaces $C+o(1)$ for all sufficiently large $s$.
\end{theorem}
The new ingredients are near-balanced matching-potential packing, a thin additive precision
gap, exact counting of crossing packets and a larger grouping alphabet at the root than in
the tail. Exact taper, sharp precision and factorial-Cauchy modes come from the supplied
reports and are not claimed as new here. Small or exceptional parameter choices are handled
by direct summation, as in the source. The crossover can be enormous.

\subsection{Deterministic near-balanced packing for multigraphs}
Let an active graph have $M$ edges and maximum degree $\Delta$.
Let $q=r^a$ be the number of labels on each side and $\dd=q^2$ the number of ordered
label cells. Greedy edge coloring produces $K\le2\Delta-1$ matchings $E_j$,
including for multigraphs: an edge has at most $2\Delta-2$ already colored neighbors.
Write $m_j=|E_j|$.

Assign independent uniform labels in $[q]$ to vertices. For each matching and label cell
$\gamma$, let $X_{j,\gamma}$ be its load. Put
\[
 z=1+\lambda^{-1},\qquad
 t_j=\left\lceil(1+\lambda^{-1})m_j/\dd+10\lambda^2\right\rceil,
 \qquad \Phi=\sum_{j,\gamma}z^{X_{j,\gamma}-t_j}.
\]
Distinct edges within a matching have independent endpoints. Independence between different
matchings is neither true nor needed, since expectation is linear. Thus
\begin{align*}
 \mathbb E\Phi
 &=\sum_{j,\gamma}z^{-t_j}(1+1/(\lambda\dd))^{m_j}\\
 &\le K\dd\,e^{-5\lambda}<1
\end{align*}
for all sufficiently large $s$ in the schedules below. To see the bound per term, use
\[
 \log(1+1/\lambda)\ge\frac1{2\lambda},\qquad
 (1+1/\lambda)\log(1+1/\lambda)\ge\frac1\lambda,
 \qquad \log(1+x)\le x.
\]
The $m_j$ terms cancel in the exponent and the additive $10\lambda^2$ leaves at most
$-5\lambda$. The schedule has $\dd\le s$ and $K\le4s$, so the final strict inequality
holds eventually. Logarithms in this expectation estimate are natural.

\paragraph{Exact derandomization and bit cost.}
Condition on a partial vertex labeling. Different edges of one matching still depend on
disjoint unassigned vertices. The probability that an edge receives cell $\gamma$ is in
$\{0,q^{-2},q^{-1},1\}$. Therefore the conditional expectation of each summand is exactly
\[
 z^{-t_j}\prod_{e\in E_j}(1+p_e/\lambda).
\]
At each vertex, try all $q$ possible labels and choose one minimizing this rational number
summed over all terms. The average over labels is the current conditional expectation,
so the invariant $\mathbb E[\Phi\mid\text{assigned labels}]<1$ persists.
At completion $X_{j,\gamma}<t_j$ for every cell and matching, since otherwise one summand
would be at least one. Hence, using $\lceil x\rceil-1<x$,
\begin{equation}\label{eq:load}
 X_\gamma<(1+\lambda^{-1})M/\dd+10K\lambda^2
 \le(1+\lambda^{-1})M/\dd+20\Delta\lambda^2.
\end{equation}
There are polynomially many vertices, tried labels, cells and edges in $s+\mathcal L$.
Every exponent $t_j$ is $O(s+\lambda^2)$, the rational bases have polynomial-size
integer numerators and denominators, and products involve at most polynomially many factors.
Exact rational products, sums and comparisons therefore have polynomial bit length and
polynomial time/space cost. This deliberately simple implementation is not claimed practical.

Pad every cell to the same multiple $H_{\rm pad}$ of packet length $L=2^a$.
The total padded size is
\begin{equation}\label{eq:padding}
 B=H_{\rm pad}\dd
 \le(1+\lambda^{-1})M+20\dd\Delta\lambda^2+L\dd.
\end{equation}
Consequently $B=(1+o(1))M$ whenever
$\dd\Delta\polylog(s)=o(M)$ and $L\dd=o(M)$.
This replaces the source's sufficient $\dd^2\Delta=o(M)$ condition while leaving its
subsequent packet arrangement intact. There are $B/L$ packets. Each coordinate-symbol
class has at most $B/r$ real ports on either side, so a child has at most
\[
 2B/r=(2/r)(1+o(1))M
\]
edges. New packet vertices have degree at most $2L$; retained vertices keep their old degrees.

\paragraph{A sharp scale for this specific packing architecture.}
Take $t$ disjoint pairs, with $\Delta$ parallel edges on each pair, so $M=t\Delta$.
Every nonzero cell load is a multiple of $\Delta$ and at least one is nonzero.
Equal padding therefore forces $B\ge\Delta\dd$. Near-unit total padding is impossible
if $\Delta\dd/M$ diverges. This is only a limitation of equal-padded collision packing,
not a lower bound for contraction or SAT.

\paragraph{Count crossing packets rather than all packets.}
At any fixed coordinate, exactly the fraction $1-1/r$ of ordered padded label pairs
have unequal symbols. Every packet fixes that pair. Thus exactly
$(1-1/r)B/L$ packets can have crossing modes there. At retained valuation $i$, put
$\Dt=2^{a-i}$. The number $c$ of column-tile families satisfies
\begin{equation}\label{eq:families}
 c\le2(1-1/r)B/\Dt=(2(1-1/r)+o(1))M/\Dt.
\end{equation}
Empty or purely local tiles only reduce this number.

\subsection{Thin precision gaps and exact taper}
Set $g(\lambda)=(\lceil\log_2\lambda\rceil)^2$. For positive integer taper lengths $u_d$, use
\begin{equation}\label{eq:thin}
 a_d=\sum_{j=0}^d u_j+(d+1)g(\lambda),\qquad a_{-1}=0.
\end{equation}
The source requires the cumulative packet exponent to dominate inherited precision;
it does not require literally doubling every taper increment.
At a retained $i<u_d$,
\[
 \Dt=2^{a_d-i}\ge2^{a_{d-1}+g+1}.
\]
Use the complete linear taper, with $a=a_d$, $u=u_d$,
\begin{equation}\label{eq:taper}
 h=\left\lceil\frac{M(2L-1)+1}{L(2u-2+2^{1-u})}\right\rceil,
 \qquad q_i=h(u-i)\quad(0\le i<u),
\end{equation}
and $q_i=0$ thereafter. The exact condition count is
\[
 \sum_{i=0}^{u-1}\frac{L}{2^i}(u-i)=L(2u-2+2^{1-u}),
\]
so the total multiplicity is greater than the deformation degree $M(2L-1)$.
The supplied sharp precision identity is
\begin{equation}\label{eq:precision}
 p_t=p+B_i+(q_i-1)(a+1-i),\qquad
 B_i=\Dt\sum_{j=1}^{a-i}j2^{-j}(q_i-q_{i+j})\le6h\Dt.
\end{equation}
The taper is nonincreasing, its precision increments are positive, and the immutable
creation-precision lifting convention is unchanged.

Along a recursion path, precision increments are at most
$s\lambda^{O(1)}2^{a_d}$. Inherited precision is therefore at most
$s\lambda^{O(1)}2^{a_{d-1}}$. With leaf threshold $H=\lceil s/\lambda\rceil$,
active calls have $M>s/\lambda$ and $u=O(\lambda)$, so inherited precision divided
by $M\Dt/u$ is at most $\lambda^{O(1)}2^{-g}=o(1)$. The last evaluation-loss term
in \eqref{eq:precision} is negligible for the same reason. If $u$ grows as a positive
power of $\lambda$ uniformly at active depths, then
$h=(1+o(1))M/u$ and
\begin{equation}\label{eq:precision6}
 p_t\le(6+o(1))M\Dt/u.
\end{equation}
The depth proof below supplies exactly that uniform growth. Over $O(\log\lambda)$ levels,
the total gap consumes $O((\log\lambda)^3)=o(\lambda)$ label coordinates.
No extra root precision is assumed.

\subsection{Factorial modes and the complete arithmetic interface}
Use the supplied exact Cauchy features, whose guaranteed weights are
\[
 w_b=1+b+2v_2(b!)=3b+1-2s_2(b),\qquad 0\le b<\Dt/2,
\]
where $s_2(b)$ is binary popcount. The actual row parameters satisfy the required
prime-subring congruence at every interpolation value: $z_x/2$ differs from a rational
2-integral quantity $2^i/d_0$ by a multiple of $2^{a+1-i}$.
Since $b<\Dt/2=2^{a-i-1}$, the factorial-divisibility argument applies.
There is no division by a factorial in a residue ring. Binomial features are computed
as ordinary integers; $b!$ is multiplied into the target coefficient. Introduced denominators
are units. Changing $r$, the taper or the packing does not change these local facts.

For $c$ families and precision $p$, the generating-product estimate is
\[
 \log N(p,c)\le\left(\frac{\pi\sqrt2}{3}+o(1)\right)\sqrt{cp}.
\]
Apply monotonicity using uniform upper envelopes
\[
 P=(6+o(1))M\Dt/u,\qquad
 C_{\rm fam}=(2(1-1/r)+o(1))M/\Dt,
\]
not an assumption that actual populated families or actual precision attain these envelopes.
Their ratio has scale $\Dt^2/u$, and $\Dt\ge2^g$, so the small-saddle regime and
its errors hold uniformly. Finite tile truncation only decreases the count. Combining
\eqref{eq:families} and \eqref{eq:precision6} gives the local natural-log coefficient
\begin{equation}\label{eq:localentropy}
 \log N\le\left[2\pi\sqrt{\tfrac23(1-1/r)}+o(1)\right]M/\sqrt u.
\end{equation}

The exact local kernel, creation-time representatives and recursive child definitions are
unchanged. The original source's packet-entry enumeration costs $\exp(O(L\log s))$,
which is already lower order for the schedule below. The supplied streamed finite-difference
replacement can reduce this to $\poly(s+\mathcal L)2^L$ using at most $L$ extra internal
precision bits, but it is optional and not needed for this theorem. If used, the numerator
must be computed modulo $2^{p_{\rm creation}+v_2(\prod\beta!)}$ before dividing canonical
coordinates by the explicit power of two. Temporary precision must not alter tensor creation precision.

The extension field still has degree $4\lambda$: deformation degrees $M(2L-1)$ are
less than $s^4$ eventually. All precisions, field coordinates, rational-recovery operands and
stored descriptions have polynomial bit length. The initial entry allowance
$2^{A_0s^{0.9}}$ is lower order than the target exponent. Old entry routines are not
recursively composed; descriptions remain immutable. The bound permits $\ell\le s$;
the scalar parity specialization $\ell=1$ is compatible but unnecessary here.

\subsection{An explicit two-stage grouping schedule}
Choose exact rationals
\begin{equation}\label{eq:schedule}
 r_0=37,\quad u_0=\left\lfloor\frac{10}{1043}\lambda\right\rfloor;
 \qquad r_d=8,\quad
 u_d=\left\lfloor\frac{\lambda}{236(5/2)^{d-1}}\right\rfloor\quad(d\ge1).
\end{equation}
Use \eqref{eq:thin} and $H=\lceil s/\lambda\rceil$.
First consider all candidate depths $d\le N=\lceil\log_3\lambda\rceil+3$,
without assuming this is the actual depth. For these candidates the total gaps are
$O((\log\lambda)^3)=o(\lambda)$, and
\[
 a_0/\lambda\le10/1043+o(1),\qquad
 a_d/\lambda\le A+o(1),\quad
 A=\frac{10}{1043}+\frac5{708}=\frac{12295}{738444}.
\]
The packing exponents are strictly less than one:
\begin{align}
 0.9+2(10/1043)\log_2 37&=0.999893640759903\ldots<1,\label{eq:rootexp}\\
 0.9+2A\log_2 8=0.9+6A&=\frac{307654}{307685}<1.\label{eq:tailexp}
\end{align}
The first follows from the exact integer inequality $37^{200}<2^{1043}$;
the second is rational. Thus $\dd\Delta\lambda^3/s=o(1)$ uniformly in the candidate
range. Also $L\dd\lambda/s=o(1)$: its root exponent is
$(10/1043)(1+2\log_2 37)$ and its tail exponent is $7A$, both below one.
Finally $2L\le\lceil s^{0.9}\rceil$ eventually, preserving the degree interface.

\paragraph{Noncircular depth bootstrap.}
For every candidate active tail call, these estimates give $B/M\le4/3$ eventually,
so each child has at most $M/3$ edges. Starting from $M\le s$, an active call cannot
reach depth $N$; it would already have at most $H$ edges. Hence every real active call
lies in the candidate range, closing the induction. This coarse bound ensures
$u_d$ grows as a positive power of $\lambda$, since $5/2<3$.
More sharply, \eqref{eq:padding} gives
$B/M=1+O(1/\lambda)+o(1/\lambda)$ uniformly. Over $O(\log\lambda)$ levels this
accumulates to $1+o(1)$, yielding for $d\ge1$
\begin{equation}\label{eq:pathsize}
 M_d\le(1+o(1))s\frac2{37}\,4^{-(d-1)}.
\end{equation}
An active tail depth is at most $\log_4\lambda+O(1)$, and therefore
\[
 u_d=\Omega(\lambda^{\eta+o(1)}),\qquad
 \eta=1-\frac{\log(5/2)}{\log4}=0.33903595\ldots>0.
\]
This justifies uniform floor errors, exact-taper asymptotics and the inherited-precision estimates.
The finite construction can check all guards and use direct summation when they fail.

\subsection{The coefficient and total resources}
Write
\[
 K(r)=\frac{2\pi}{\ln2}\sqrt{\frac23(1-1/r)}.
\]
Convert \eqref{eq:localentropy} to base two and sum along \eqref{eq:pathsize}. The geometric
ratio in the tail is $\sqrt{5/2}/4<1$. Thus
\begin{align}
 C&=\frac{K(37)}{\sqrt{10/1043}}
 +\frac{(2/37)K(8)}{\sqrt{1/236}\,(1-\sqrt{5/2}/4)}\label{eq:C}\\
 &=84.06627149423657\ldots<85.\nonumber
\end{align}
An elementary rational upper certificate uses
\[
 \pi<22/7,\quad \ln2>69/100,\quad
 \sqrt{12516/185}<8226/1000,
\]
\[
 \sqrt{413/3}<11734/1000,\qquad \sqrt{5/2}<1582/1000.
\]
These imply
\[
 C<\frac{9128012476}{108030195}=84.49501064\ldots<85.
\]
The logarithm lower bound also follows by taking the first three positive terms of
$\ln2=2\sum_{j\ge0}(1/3)^{2j+1}/(2j+1)$.
An independent checker uses slightly different rational envelopes and obtains another
strict upper bound $1659638632/19727253<85$.

Polynomial branching overhead over $O(\log\lambda)$ levels contributes
$O(\lambda\log\lambda)$ to the exponent, while leaf work contributes $O(s/\lambda)$.
Both are lower order. So is original packet enumeration, since $L=s^{A+o(1)}$ with
$A<0.02<0.9$; in particular this entry cost also fits the original $2^{O(s^{0.9})}$ allowance. Depth-first accumulation stores only polynomially many bits. Branching depends on
$s$, not on $\mathcal L$, so the dependence on description length remains a fixed polynomial.
This completes the relative proof of Theorem~\ref{thm:85}.

\subsection{What the numerical comparison does and does not mean}
\begin{center}
\begin{tabular}{@{}lr@{}}\toprule
Declared normalized construction & Leading certificate\\\midrule
Supplied contraction schedule & $648.4410666487$\\
That schedule plus factorial weights alone & $374.3776\ldots$\\
Matching packing with one fixed alphabet & $93.8049738133$\\
This two-stage schedule & $84.0662714942$\\
Conservative eventual certificate here & $85$\\\bottomrule
\end{tabular}
\end{center}
These are comparable certificates, not measured speeds or lower bounds on earlier algorithms.
The $648.4411$ construction uses geometric modes, $r=256$, threefold padding and doubled
cumulative taper; the naive factorial substitution divides its leading coefficient by
$\sqrt3$ after compatibility checks. A single-stage matching construction with $r=46$,
taper decay $8$ and initial taper fraction $10/1263$ yields the intermediate value shown.
The two-stage schedule is not claimed globally optimal.

All these bounds still have denominator $\sqrt{\log s}$. The unrestricted SAT wire budget
remains $O_k(n\log\log n)$ and the fixed-density budget $O_{k,D}(n)$.
Optimal finite taper profiles are compatible but have the same leading six-bit loss;
they do not supply a new logarithmic power. A true order improvement needs a different
summation or representation idea, not merely further tuning of these constants.
